\subsection{Three further lifts and the exact counts}\label{sec:searched-lifts}
The same formulas produce smaller counts. Let $d$ be a polynomial in $D$ and put
\[
 \delta_1^+=\delta_2^+=d,\qquad G_i^+=a_iDd,\qquad S^+=bDd^2.
\]
The three further choices are
\begin{center}
\begin{tabular}{cclll}
\toprule
Choice & $d$ & $(a_1,a_2,b)$ & $\lambda_1^+$ & $\lambda_2^+$\\
\midrule
A & $D(D-3)$ & $(10,15,30)$ & $4D^2-15D+12$ & $3D^2-10D+4$\\
B & $D(D-3)$ & $(14,7,14)$ & $D$ & $3D^2-10D+4$\\
C & $(D-1)(D-2)$ & $(14,7,14)$ & $D$ & $3D^2-10D+4$\\
\bottomrule
\end{tabular}
\end{center}
Their remaining bottom factors are obtained without matrix inversion:
\[
 W_0^+=\frac{a_1a_2}{b}D,\quad
 U_i^+=\frac{a_i(a_i-1)}bD,\quad
 V_i^+=\frac{a_i}{b}(a_iD^2-\lambda_i^+).
\]
For A these are $5D,3D,7D,2D^2+5D-4,6D^2+5D-2$, in the order $W_0,U_1,U_2,V_1,V_2$. For B and C they are $7D,13D,3D,14D^2-D,2D^2+5D-2$.

\begin{proof}[Proof of Theorem~\ref{thm:lifts-main}]
Every displayed basis is nonsingular, so its table is associative. In A the new first cell factor has Laurent sequence $(4,1,6,1,4)$, which dominates that of $D$ with the required parity. Its bottom factor differs from $V_1=5D-4$ by $2D^2$. All other effectivity checks are those already proved for the first lift.

For B the first cell factor equals the trace factor $D$, while the second is unchanged. The new bottom-factor differences for $W_0,U_1,U_2,V_1,V_2$ are respectively
\[
 6D,\quad12D,\quad2D,\quad14D^2-6D+4,\quad10D-6.
\]
All have nonnegative even Laurent coefficients. The differences $G_1^+-G_1$ and $G_2^+-G_2$, on weights $[-3,3]$, are
\[
 (14,38,30,12,30,38,14),\qquad(6,16,10,0,10,16,6).
\]
The difference $S^+-S$, on $[-5,5]$, is
\[
 (14,52,60,16,6,24,6,16,60,52,14).
\]
These give every remaining inequality and parity check. In C, increasing $d$ by $2$ increases $G_i^+$ and $S^+$ coefficientwise by even nonnegative polynomials; the previous inequalities therefore persist. Proposition~\ref{prop:criterion} proves that all three are effective associative lifts.

For C the factors $D,d,W_0^+,U_i^+,V_i^+$ are symmetric and unimodal. In particular $V_1^+$ and $V_2^+$ have sequences $(14,55,82,55,14)$ and $(2,13,20,13,2)$. The product $d\lambda_1^+=dD$ is unimodal, and $d\lambda_2^+$ has the already displayed sequence $(3,5,10,11,14,11,10,5,3)$. The convolution argument above proves unimodality of all $132$ coefficients.

The following table provides the counts by multiplication type. Each entry is the ordered pair (cancelling multiplicities, supported zero-trace weights). ``Cell output'' excludes the bottom output of a same-cell product. Products in both orders are included. The two earlier choices are denoted I and II, with II the unimodal one.
\begin{center}
\small
\begin{tabular}{lrrrrr}
\toprule
Type & I & II & A & B & C\\
\midrule
$\Sigma^2$ & $(672,2)$ & $(1872,2)$ & $(672,2)$ & $(160,2)$ & $(720,2)$\\
$\Sigma$ times a cell & $(3040,16)$ & $(5040,16)$ & $(3040,16)$ & $(2400,16)$ & $(4080,16)$\\
Same-cell output & $(1060,20)$ & $(1672,24)$ & $(724,20)$ & $(388,4)$ & $(664,8)$\\
Same-cell bottom & $(1580,16)$ & $(1580,16)$ & $(1620,16)$ & $(2940,16)$ & $(2940,16)$\\
Cross-cell bottom & $(1600,0)$ & $(1600,0)$ & $(1600,0)$ & $(2400,0)$ & $(2400,0)$\\
\midrule
Total & $(7952,54)$ & $(11764,58)$ & $(7656,54)$ & $(8288,38)$ & $(10804,42)$\\
\bottomrule
\end{tabular}
\end{center}
Each entry is obtained from Lemma~\ref{lem:table} by expanding $D=q+2+q^{-1}$ and summing $(d_j-|c_j|)/2$, while counting a supported zero once when $c_j=0<d_j$. Unit products contribute $(0,0)$. The displayed multiplication rule gives every input to this finite calculation. For example $S$ has weights $[-4,4]$ whereas $S^+$ has weights $[-5,5]$, accounting for the two supported zero weights in the first row. The independently checked matrix-coordinate computation and a separate word-reduction computation agree on all five tables and these counts. This proves the theorem.
\end{proof}

These constructions were found in a finite parameter search, whose bounds and complete output are recorded in the workbench. The verification here establishes their associativity, effectivity, counts, and stated unimodality. It does not turn a minimum in a search window into a minimum over arbitrary dimension rings or over all inverse-semiring lifts.
